This manual explains how to build and use the native emulator. You do not need to know its C implementation, historical Pascal source, or CPU internals to run the supplied RC4000 system.

Development and testing use Linux on x86-64. The runtime uses C17, POSIX interfaces and GNU Readline. Other hosts, including macOS, have not been validated; platform-specific setup and testing are left to interested contributors.

For supported-system boundaries and known differences from hardware and the DOS emulator, see \textit{Compatibility and Limitations} (\path{doc/COMPATIBILITY.md}).

The local terminal can address three command interpreters: the simulator, the RC system process \passthrough{\lstinline!s!}, and FP inside a user process. Commands for one must not be entered at another.

For the separate automatic TEM/PRIMO/SOS environment, read
\manualref{mips-ts}{MIPS/TS: a terminal session}. It has its own disk pair,
demonstration login and resource limits. The default-profile process-creation
and host-import exercises are not additional steps in a MIPS/TS login.

\section*{Reading commands and responses}

Examples throughout the manual have context titles: host shell,
simulator prompt, guest system process \texttt{s}, guest FP, a tool's own
input, or configuration/file contents. These titles
are not text to enter. Output examples are labelled separately; do not type
the monitor's messages or prompts. Return submits each command line unless
the accompanying instructions say otherwise.

Angle-bracketed key names such as \texttt{<Control-E>} describe a key press,
not literal input. Likewise, \texttt{local:} and \texttt{remote:} in
two-terminal transcripts identify the terminal; omit these prefixes when
typing. Ellipses in displayed results stand for session-dependent fields,
not text to enter. Monospaced names in prose identify commands, files or
configuration settings; their context still determines where to use them.

Local emulator-escape examples use Control-E, the shipped default. Substitute
your configured \texttt{OPERATOR ESCAPE} key if you change it; Control-E at
Readline's operator prompt is instead a line-editing binding. Remote Telnet
attention uses its endpoint's own setting, not the local emulator escape key.
See \manualref{keyboard-and-date-settings}{keyboard settings} and
\manualref{request-guest-attention-not-emulator-suspension}{remote attention}.

Syntax-reference blocks are not scripts: brackets indicate optional arguments,
vertical bars indicate alternatives, and uppercase placeholders name values
you must supply. Fresh-session examples are independent unless their text
explicitly continues an earlier procedure. A command returning without an
error is evidence for that example, not every function of the historical tool.

Cross-references within this manual include a page number for use with a printed
manual. References to \path{doc/} files are repository-relative paths, not
runtime files or files inside the guest disk image.

\section*{Terms used in this manual}
\begin{description}
\item[Simulator] The emulator program running on the host. Its \texttt{sim4000>} prompt controls
execution, media and debugging; it is outside the guest operating system.
\item[System process \texttt{s}] The guest supervisor's command interface for
creating and controlling processes.
\item[FP] The file processor: a guest command interpreter inside an ordinary
user process. It starts utilities, compilers and other programs.
\item[Catalog entry and scope] A guest name and its metadata. Temporary scope
belongs to a process's work; retained user scope can survive process removal.
Neither scope by itself determines whether host disk writes survive exit.
\item[Segment and RC byte] Guest allocation and addressing units.
A backing-store segment holds 256 24-bit words (768 host bytes).
An RC byte is a 12-bit halfword, not an eight-bit host byte.
\item[Overlay] Changes kept separately from the base disk image.
A volatile overlay disappears on emulator exit; a persistent overlay is
saved in its own host file and must stay with the matching base image.
\item[Attention and suspension] Guest attention selects which process receives
terminal input. Emulator suspension stops CPU execution and opens the
simulator prompt. They are not interchangeable.
\end{description}
